\section{Preliminaries and Problem Formulation}
\label{sec:problem_formulation}

This section formalizes RAG from the perspective of claim-level evidence coverage. We first define the task setting and the basic objects used throughout the paper, then introduce atomic claims, evidence support relations, coverage maps, and the budgeted optimization objective that motivates \ragname.
% 中文：本节从 claim-level evidence coverage 的角度形式化检索增强生成问题。我们首先定义任务设置和全文使用的基本对象，然后引入原子事实声明、证据支持关系、coverage map，以及支撑 \ragname 的预算约束优化目标。

\subsection{Task Setting}
\label{subsec:task_setting}

Let $q$ denote a user question and let $\mathcal{D}$ denote an external data repository, such as a document collection or a structured knowledge base. We use a unified evidence universe
\begin{equation}
    \mathcal{E} = \{e_1, e_2, \ldots, e_N\}
\end{equation}
to denote the set of retrievable evidence units derived from $\mathcal{D}$. An evidence unit $e_j$ can be a passage chunk, a database record, or any other retrievable unit that can be cited or inspected.
% 中文：令 $q$ 表示用户问题，$\mathcal{D}$ 表示外部数据仓库，例如文档集合或结构化知识库。我们用统一的证据空间 $\mathcal{E}$ 表示从 $\mathcal{D}$ 中派生出来的可检索证据单元集合。一个证据单元 $e_j$ 可以是段落块、数据库记录，或任何可以被引用和检查的可检索单元。

A retrieval-augmented generation system has access to a retriever $R$, a generator $G$, and optionally one or more verifiers $V$. Given $q$, the retriever selects a subset of evidence units
\begin{equation}
    \mathcal{E}_q = R(q, \mathcal{E}; B_R),
\end{equation}
where $B_R$ is a retrieval budget controlling the number of retrieval rounds, candidate evidence units, or retrieved tokens. The generator then produces an answer $a = G(q, \mathcal{E}_q)$, optionally with citations. In conventional RAG, the primary control target is whether $\mathcal{E}_q$ is relevant to $q$. In this paper, we instead ask whether the factual claims in $a$ are covered by evidence.
% 中文：一个检索增强生成系统可以访问检索器 $R$、生成器 $G$，以及可选的一个或多个验证器 $V$。给定问题 $q$，检索器从证据空间中选择一组证据 $\mathcal{E}_q$，其中 $B_R$ 是检索预算，可以控制检索轮数、候选证据数量或检索 token 数。随后，生成器基于 $q$ 和 $\mathcal{E}_q$ 生成答案 $a$，答案可以带有引用。传统 RAG 的主要控制目标是 $\mathcal{E}_q$ 是否与 $q$ 相关；本文关注的则是 $a$ 中的事实声明是否被证据覆盖。

The desired output of our setting is not only an answer string. Instead, the system should return
\begin{equation}
    \mathcal{O} = (a, \mathcal{C}, \mathcal{S}, \mathcal{Z}, \Gamma, \mathbf{r}),
\end{equation}
where $a$ is the final answer, $\mathcal{C}$ is the set of atomic factual claims in $a$, $\mathcal{S}$ assigns evidence sets to claims, $\mathcal{Z}$ assigns support labels to claims, $\Gamma$ is the citation assignment, and $\mathbf{r}$ contains confidence or reliability scores. This output structure makes evidence use explicit and auditable.
% 中文：在我们的任务设置中，系统的期望输出不只是一个答案字符串，而是一个结构化对象 $\mathcal{O}$。其中 $a$ 是最终答案，$\mathcal{C}$ 是答案中的原子事实声明集合，$\mathcal{S}$ 为每个 claim 分配证据集合，$\mathcal{Z}$ 为每个 claim 分配支持标签，$\Gamma$ 是引用分配，$\mathbf{r}$ 包含置信度或可靠性分数。这样的输出结构使证据使用过程显式且可审计。

\subsection{Atomic Factual Claims}
\label{subsec:atomic_claims}

An answer $a$ is decomposed into a set of atomic factual claims
\begin{equation}
    \mathcal{C}(a) = \{c_1, c_2, \ldots, c_m\}.
\end{equation}
An atomic claim is the smallest factual statement that can be independently verified against evidence. For example, the sentence ``Demis Hassabis co-founded DeepMind in London in 2010 and studied at the University of Cambridge'' can be decomposed into four atomic claims: Demis Hassabis co-founded DeepMind; DeepMind was founded in London; DeepMind was founded in 2010; and Demis Hassabis studied at the University of Cambridge.
% 中文：答案 $a$ 被拆解为一组原子事实声明 $\mathcal{C}(a)$。原子 claim 是能够独立地与证据进行验证的最小事实陈述。例如句子 “Demis Hassabis co-founded DeepMind in London in 2010 and studied at the University of Cambridge” 可以拆解为四个原子 claim：Demis Hassabis 共同创立了 DeepMind；DeepMind 创立于伦敦；DeepMind 创立于 2010 年；Demis Hassabis 就读于 University of Cambridge。

Not all claims are equally important for answering a question. We associate each claim $c_i$ with a non-negative importance weight $w_i \geq 0$. Answer-bearing claims, such as entities, dates, locations, numbers, and relations directly requested by $q$, should receive larger weights. Auxiliary explanatory or background claims may receive smaller weights. Let
\begin{equation}
    \mathbf{w} = (w_1, w_2, \ldots, w_m)
\end{equation}
denote the claim-importance vector. These weights allow the formulation to distinguish missing evidence for a core answer from missing evidence for peripheral explanation.
% 中文：并非所有 claim 对回答问题都同等重要。我们为每个 claim $c_i$ 关联一个非负重要性权重 $w_i$。问题直接询问的实体、日期、地点、数值和关系等 answer-bearing claims 应获得更高权重；辅助解释或背景性 claims 可以获得较低权重。claim-importance vector $\mathbf{w}$ 使得形式化能够区分核心答案缺证据与外围解释缺证据。

\subsection{Evidence Sets and Support Labels}
\label{subsec:evidence_support}

For each claim $c_i$, the system assigns a candidate evidence set
\begin{equation}
    S_i \subseteq \mathcal{E},
\end{equation}
where $S_i$ may contain one or multiple evidence units. We use $\mathcal{S} = \{S_i\}_{i=1}^{m}$ to denote the collection of claim-specific evidence sets. A verifier $V$ evaluates the relation between $c_i$ and $S_i$ and returns a support label
\begin{equation}
    z_i = V(c_i, S_i) \in \mathcal{Y},
\end{equation}
where
\begin{equation}
    \mathcal{Y} = \{\mathsf{SUP}, \mathsf{MSUP}, \mathsf{CONTRA}, \mathsf{NEI}\}.
\end{equation}
Here, $\mathsf{SUP}$ means that a single evidence unit directly supports the claim; $\mathsf{MSUP}$ means that multiple evidence units jointly support the claim; $\mathsf{CONTRA}$ means that the evidence contradicts the claim; and $\mathsf{NEI}$ means not enough information.
% 中文：对于每个 claim $c_i$，系统分配一个候选证据集合 $S_i$，其中可以包含一个或多个证据单元。我们用 $\mathcal{S}$ 表示所有 claim-specific evidence sets 的集合。验证器 $V$ 评估 $c_i$ 与 $S_i$ 之间的关系，并返回支持标签 $z_i$。其中 $\mathsf{SUP}$ 表示单个证据单元直接支持该 claim；$\mathsf{MSUP}$ 表示多个证据单元共同支持该 claim；$\mathsf{CONTRA}$ 表示证据与 claim 矛盾；$\mathsf{NEI}$ 表示证据不足。

The verifier may also return a confidence score $r_i \in [0,1]$, producing
\begin{equation}
    (z_i, r_i) = V(c_i, S_i).
\end{equation}
In practice, $V$ can be implemented by an ensemble of natural language inference models, question answering models, rule-based consistency checks, and LLM-based judges. The formulation does not assume that the verifier is perfect. Instead, verification is treated as an explicit and inspectable approximation to the latent evidence-support relation.
% 中文：验证器还可以返回置信度分数 $r_i$。在实际实现中，$V$ 可以由自然语言推理模型、问答模型、规则一致性检查和 LLM judge 组成的 ensemble 来实现。这里的形式化并不假设验证器是完美的，而是将验证视为对潜在 evidence-support relation 的显式、可检查近似。

We define a binary support indicator
\begin{equation}
    s(c_i, S_i) =
    \mathbb{I}\left[z_i \in \{\mathsf{SUP}, \mathsf{MSUP}\}\right],
    \label{eq:support_indicator}
\end{equation}
where $\mathbb{I}[\cdot]$ is the indicator function. A contradicted or insufficiently evidenced claim is therefore not considered covered, even if the evidence set is topically relevant.
% 中文：我们定义二值支持指示函数 $s(c_i,S_i)$。当 $z_i$ 为 $\mathsf{SUP}$ 或 $\mathsf{MSUP}$ 时，该 claim 被视为 covered；当 claim 被反驳或证据不足时，即使证据集合在主题上相关，也不视为被覆盖。

\subsection{Claim-Evidence Coverage Map}
\label{subsec:coverage_map}

The central intermediate object in our formulation is a claim-evidence coverage map. We define it as
\begin{equation}
    \mathcal{M} = (\mathcal{C}, \mathcal{E}^{\star}, \mathcal{A}, \mathcal{Z}, \mathbf{r}),
    \label{eq:coverage_map}
\end{equation}
where $\mathcal{E}^{\star} \subseteq \mathcal{E}$ is the evidence used by the system, $\mathcal{A}$ is a set of support edges or hyperedges connecting claims to evidence units, $\mathcal{Z}$ is the set of support labels, and $\mathbf{r}$ contains confidence scores. For direct support, an edge may connect one claim to one evidence unit, $(c_i, e_j)$. For multi-hop support, a hyperedge may connect one claim to a set of evidence units, $(c_i, \{e_{j_1}, \ldots, e_{j_k}\})$.
% 中文：本文形式化中的核心中间对象是 claim-evidence coverage map。我们将其定义为 $\mathcal{M}$，其中 $\mathcal{E}^{\star}$ 是系统实际使用的证据子集，$\mathcal{A}$ 是连接 claim 与证据单元的支持边或超边，$\mathcal{Z}$ 是支持标签集合，$\mathbf{r}$ 是置信度分数集合。对于直接支持，一条边可以连接一个 claim 和一个证据单元；对于多跳支持，一条超边可以连接一个 claim 和多个证据单元。

The coverage map differs from ordinary citation metadata. A citation only records that a span of generated text points to a source, whereas $\mathcal{M}$ records whether the source actually supports each atomic claim. Thus, $\mathcal{M}$ can represent supported claims, multi-hop supported claims, contradicted claims, unsupported claims, missing evidence, and citation mismatches in one structure.
% 中文：coverage map 不同于普通 citation metadata。引用只记录生成文本的某个 span 指向某个来源，而 $\mathcal{M}$ 记录该来源是否真正支持每个原子 claim。因此，$\mathcal{M}$ 可以在同一结构中表示被支持的 claims、多跳支持的 claims、被反驳的 claims、未被支持的 claims、缺失证据以及引用不匹配。

Let $\Gamma(a)$ denote the citation assignment in the final answer, where $\Gamma(c_i) \subseteq S_i$ gives the evidence units cited for claim $c_i$. Citation faithfulness requires not only that citations exist, but that cited evidence entails the corresponding claim:
\begin{equation}
    \forall c_i \in \mathcal{C}(a), \quad
    \Gamma(c_i) \neq \emptyset
    \Rightarrow
    s(c_i, \Gamma(c_i)) = 1.
    \label{eq:citation_faithfulness_condition}
\end{equation}
This condition highlights the difference between citation presence and citation support.
% 中文：令 $\Gamma(a)$ 表示最终答案中的引用分配，其中 $\Gamma(c_i)$ 给出 claim $c_i$ 所引用的证据单元。citation faithfulness 不仅要求存在引用，还要求被引用证据能够蕴含对应 claim。这个条件强调了“有引用”和“引用真正支持 claim”之间的区别。

\subsection{Coverage Metrics}
\label{subsec:coverage_metrics}

Given a coverage map $\mathcal{M}$, the unweighted claim coverage is
\begin{equation}
    \mathrm{Cov}(\mathcal{C}, \mathcal{S})
    =
    \frac{1}{|\mathcal{C}|}
    \sum_{i=1}^{|\mathcal{C}|}
    s(c_i, S_i).
    \label{eq:claim_coverage}
\end{equation}
Because answer-bearing claims are more important than background claims, we also define weighted evidence coverage:
\begin{equation}
    \mathrm{WCov}(\mathcal{C}, \mathcal{S}, \mathbf{w})
    =
    \frac{
        \sum_{i=1}^{|\mathcal{C}|} w_i \cdot s(c_i, S_i)
    }{
        \sum_{i=1}^{|\mathcal{C}|} w_i + \epsilon
    },
    \label{eq:weighted_coverage}
\end{equation}
where $\epsilon$ is a small constant for numerical stability. $\mathrm{WCov}$ gives higher penalty to missing evidence for core answer claims.
% 中文：给定 coverage map，普通 claim coverage 是被支持 claims 所占比例。由于 answer-bearing claims 比背景 claims 更重要，我们进一步定义加权证据覆盖率 $\mathrm{WCov}$。其中 $w_i$ 是 claim 重要性权重，$\epsilon$ 是用于数值稳定的小常数。$\mathrm{WCov}$ 会对核心答案 claim 的缺证据情况施加更高惩罚。

We further define the unsupported claim rate and contradiction rate as
\begin{equation}
    \mathrm{Unsup}(\mathcal{C}, \mathcal{S})
    =
    \frac{1}{|\mathcal{C}|}
    \sum_{i=1}^{|\mathcal{C}|}
    \mathbb{I}\left[z_i = \mathsf{NEI}\right],
    \label{eq:unsupported_rate}
\end{equation}
and
\begin{equation}
    \mathrm{Contra}(\mathcal{C}, \mathcal{S})
    =
    \frac{1}{|\mathcal{C}|}
    \sum_{i=1}^{|\mathcal{C}|}
    \mathbb{I}\left[z_i = \mathsf{CONTRA}\right].
    \label{eq:contradiction_rate}
\end{equation}
These metrics are necessary because an answer may achieve high lexical overlap or answer F1 while still containing unsupported or contradicted factual claims.
% 中文：我们进一步定义 unsupported claim rate 和 contradiction rate。前者表示证据不足的 claims 比例，后者表示被证据反驳的 claims 比例。这些指标是必要的，因为一个答案即使具有较高的词面重合度或 QA-F1，也仍然可能包含未被支持或被证据反驳的事实声明。

Finally, citation faithfulness can be measured over cited claims as
\begin{equation}
    \mathrm{CFaith}(a, \Gamma)
    =
    \frac{
        \sum_{c_i \in \mathcal{C}_{\Gamma}}
        \mathbb{I}\left[s(c_i, \Gamma(c_i)) = 1\right]
    }{
        |\mathcal{C}_{\Gamma}| + \epsilon
    },
    \label{eq:citation_faithfulness}
\end{equation}
where $\mathcal{C}_{\Gamma} = \{c_i \in \mathcal{C}(a): \Gamma(c_i) \neq \emptyset\}$ is the set of claims with citations. This metric evaluates whether citations support the claims they are attached to, not merely whether citations are present.
% 中文：最后，我们可以在带引用的 claims 上定义 citation faithfulness。$\mathcal{C}_{\Gamma}$ 是具有引用的 claims 集合。该指标考察引用是否真正支持其所附着的 claim，而不仅仅考察答案中是否出现了引用符号。

\subsection{Budgeted Evidence Control}
\label{subsec:budgeted_control}

Verification and repeated retrieval can improve faithfulness, but they also introduce cost and latency. We therefore define a budget vector
\begin{equation}
    \mathbf{B} = (B_R, B_E, B_V, B_T, B_L),
\end{equation}
where $B_R$ limits retrieval rounds, $B_E$ limits the number of evidence units, $B_V$ limits verification calls, $B_T$ limits generated or processed tokens, and $B_L$ limits latency. A RAG controller $\pi$ must produce an output $\mathcal{O}$ while satisfying
\begin{equation}
    \mathrm{Cost}(\pi; q, \mathcal{E}, G, V) \leq \mathbf{B}.
    \label{eq:budget_constraint}
\end{equation}
The inequality is understood component-wise when multiple budget dimensions are considered.
% 中文：验证和重复检索可以提升忠实性，但也会引入成本和延迟。因此，我们定义预算向量 $\mathbf{B}$，其中 $B_R$ 限制检索轮数，$B_E$ 限制证据单元数量，$B_V$ 限制验证调用次数，$B_T$ 限制生成或处理 token 数，$B_L$ 限制延迟。RAG 控制器 $\pi$ 必须在满足预算约束的前提下产生输出 $\mathcal{O}$。当存在多个预算维度时，该不等式按分量理解。

This budgeted view is important. The objective of \ragname is not to retrieve repeatedly for every unsupported statement, nor to make every answer cheaper than one-shot RAG. Instead, the objective is to allocate verification and retrieval budget to claims where additional evidence can most improve answer correctness and faithfulness. Supported claims should be preserved, low-value unsupported claims may be pruned or rewritten, and critical unsupported claims should trigger targeted retrieval.
% 中文：这种预算视角非常重要。\ragname 的目标不是对每个 unsupported statement 都反复检索，也不是保证每个答案都比一次性 RAG 更便宜。相反，它的目标是将验证和检索预算分配给那些额外证据最可能提升答案正确性和忠实性的 claims。已经被支持的 claims 应被保留，低价值 unsupported claims 可以被删除或改写，而关键 unsupported claims 应触发定向检索。

\subsection{Optimization Objective}
\label{subsec:optimization_objective}

We now define the claim-level evidence coverage problem. Given a question $q$, evidence universe $\mathcal{E}$, generator $G$, verifier $V$, and budget vector $\mathbf{B}$, the goal is to learn or design a controller $\pi$ that returns an answer and coverage map
\begin{equation}
    (a, \mathcal{M}) = \pi(q, \mathcal{E}, G, V; \mathbf{B})
\end{equation}
that maximize a utility function combining answer correctness, evidence coverage, citation faithfulness, and cost:
\begin{align}
    \max_{\pi} \quad
    \mathbb{E}_{q \sim \mathcal{Q}}
    \Big[
        &\mathrm{Correct}(a, q)
        + \lambda_1 \mathrm{WCov}(\mathcal{C}, \mathcal{S}, \mathbf{w}) \nonumber \\
        &+ \lambda_2 \mathrm{CFaith}(a, \Gamma)
        - \lambda_3 \mathrm{Unsup}(\mathcal{C}, \mathcal{S}) \nonumber \\
        &- \lambda_4 \mathrm{Contra}(\mathcal{C}, \mathcal{S})
        - \lambda_5 \mathrm{Cost}(\pi)
    \Big] \nonumber \\
    \text{s.t.} \quad
        &\mathrm{Cost}(\pi; q, \mathcal{E}, G, V) \leq \mathbf{B}.
    \label{eq:optimization}
\end{align}
Here, $\lambda_1,\ldots,\lambda_5$ are non-negative trade-off coefficients, and $\mathcal{Q}$ denotes the distribution of user questions.
% 中文：现在我们定义 claim-level evidence coverage problem。给定问题 $q$、证据空间 $\mathcal{E}$、生成器 $G$、验证器 $V$ 和预算向量 $\mathbf{B}$，目标是学习或设计一个控制器 $\pi$，输出答案和 coverage map，并最大化由答案正确性、证据覆盖率、引用忠实性和成本共同组成的效用函数。其中 $\lambda_1,\ldots,\lambda_5$ 是非负权衡系数，$\mathcal{Q}$ 表示用户问题分布。

\autoref{eq:optimization} should be interpreted as a conceptual objective rather than a fully differentiable training loss. At inference time, the true correctness of $a$ is generally unknown, and the support relation between claims and evidence must be approximated by retrieval, verification, and citation checking. \ragname is therefore a budgeted controller that approximates this objective by explicitly constructing and updating $\mathcal{M}$ during generation.
% 中文：公式~\ref{eq:optimization} 应被理解为概念性目标，而不是一个完全可微的训练损失。在推理时，答案 $a$ 的真实正确性通常是未知的，claims 与证据之间的支持关系必须通过检索、验证和引用检查来近似。因此，\ragname 是一个预算约束控制器，它通过在生成过程中显式构建和更新 $\mathcal{M}$ 来近似这一目标。

\subsection{Problem Statement}
\label{subsec:problem_statement}

\textbf{Problem.} Given a question $q$, an evidence universe $\mathcal{E}$, a generator $G$, a verifier $V$, and a budget vector $\mathbf{B}$, produce an answer $a$ and a claim-evidence coverage map $\mathcal{M}$ such that: (i) answer-bearing claims in $a$ are maximally supported by direct or multi-hop evidence; (ii) unsupported or contradicted claims are corrected, removed, or marked as insufficiently evidenced; (iii) citations point to evidence units that actually support the corresponding claims; and (iv) all retrieval, verification, and generation operations remain within the budget $\mathbf{B}$.
% 中文：问题定义：给定问题 $q$、证据空间 $\mathcal{E}$、生成器 $G$、验证器 $V$ 和预算向量 $\mathbf{B}$，系统需要产生答案 $a$ 和 claim-evidence coverage map $\mathcal{M}$，使得：第一，答案中的关键 answer-bearing claims 尽可能被直接证据或多跳证据支持；第二，未被支持或被反驳的 claims 被修正、删除，或标注为证据不足；第三，引用指向能够真正支持对应 claim 的证据单元；第四，所有检索、验证和生成操作都满足预算约束。

This formulation makes the gap between document relevance and answer faithfulness explicit. A set of retrieved documents can be highly relevant to $q$ while still failing to cover one or more answer claims. Conversely, an answer can be concise and well-cited only when the system controls generation at the level of claims and evidence, rather than at the level of documents alone.
% 中文：这一形式化显式揭示了文档相关性与答案忠实性之间的差距。一组检索文档可以与问题高度相关，却仍然无法覆盖一个或多个答案 claim。反过来，只有当系统在 claim 与 evidence 的层面控制生成，而不只是停留在文档层面时，答案才可能既简洁又具有可信引用。
